\section{Physical skeleton morphism and uniform complementary inverse}
\label{sec:uniforminverse}

This section uses the complete fixed-point and physical-resolvent proof
retained in \nolinkurl{references/LANE_A_ELIMINATION_INPUTS.md},
Sections 1, 10, 11, 12 and 15. The input concerns the identical open
box, every original link and face, and the actual positive unit vacuum.
Here is its exact dictionary. Its $g_{\rm YM}$ is our $g$, its
$\mathcal E$ is our $E_0$, and its coordinate action
$\alpha_h(U)_e=h_{s(e)}U_eh_{t(e)}^{-1}$ obeys
\[
 (h\cdot U)_e=h_{s(e)}^{-1}U_eh_{t(e)}
                  =(\alpha_{h^{-1}}U)_e.
\]
Thus its Hilbert-space representation
$\mathsf T_hu=u\circ\alpha_{h^{-1}}$ is exactly
$\mathsf T_hu(U)=u(h\cdot U)$ here. The identity map on functions
intertwines these representations, on $L^2$, $H^1$ and $H^2$;
no assertion that parameter inversion is a group homomorphism is needed.
The two Hamiltonians and their domains coincide. Uniqueness of the
positive unit vacuum then proves equality of the two $\psi$, of
$\rho=\psi^2$, and of the ground energies. The source scalar in its
(12.7) will be written $c_{\rm vac}=\langle1,\psi\rangle$ here;
it is not the density constant $\alpha$ in \eqref{eq:finitegap}.

For precision, the proved input is
\begin{equation}\label{eq:stronggap}
 \delta_{\rm G}\geq\Delta_{\rm sc}
 :=3\kappa\left(1-\frac{512}{3}\xi\right)
 =\kappa(3-512\xi)\geq\frac{287\kappa}{96},
 \qquad 0<\xi\leq\frac1{49152}.
\end{equation}
Its entire proof is retained, not a proposed additional hypothesis.
The exact source construction maps a collection of nonconstant
Haar sectors $z_I$ to $\phi=e^{\mathfrak C}1$, where
$\mathfrak C=\sum_{I\ne\varnothing}|z_I\rangle\langle1_I|
\otimes I_{I^c}$ is bounded and nilpotent in each finite box.
The inverse is $e^{-\mathfrak C}$, with both exponentials finite
polynomials preserving $D(H_0)=H^2$. Its full fixed-point proof gives
\[
 \psi=c_{\rm vac}\phi,\quad
 c_{\rm vac}>0,\quad c_{\rm vac}^2\|\phi\|^2=1,\quad
 e^{-\mathfrak C}(H-E_0)e^{\mathfrak C}=H_0+\mathcal K,
 \qquad H_0=\kappa\sum_eE_e.
\]
These maps preserve the vacuum scalar and all of $2bM$ in the
eigenvalue equation. The source proves that every iterate, hence both
similarities, commutes with every vertex gauge transformation.
In the orthogonal sector decomposition a physical nonempty support
has no degree-one vertex: averaging at such a vertex would give both
the same invariant vector and its zero one-link Haar average.
Every nonempty surviving support therefore contains at least four
edges of this open bipartite cubic graph. With the original
$E_e$ Casimir, $H_0\geq3\kappa$ on the physical nonconstant sectors.
The complete interaction estimate proved there is
$\|\mathcal K f\|_{\oplus,1}\leq(512\xi/3)\|H_0f\|_{\oplus,1}$,
where $\|f\|_{\oplus,1}=\sum_I\|\Pi_If\|$ and
$\|f\|\leq\|f\|_{\oplus,1}\leq2^{N/2}\|f\|$.
The retained proof supplies every commutator and counting estimate
and the contraction on $\|z\|_*\leq1/16$, with constant at most
$11/36$ in the displayed range. The sector norm equivalence is
explicitly volume dependent. The physical Neumann resolvent excludes
every $0<t<3\kappa(1-512\xi/3)$, transfers under the displayed
domain-preserving similarity, and proves \eqref{eq:stronggap} for the
original self-adjoint operator. No form positivity of a nonunitary
similarity or deletion of its norm factors is used.

The retained proof derives numerical constants for the creation method
of Yarotsky~\cite{Yarotsky}, Section 2, source equations
\texttt{mmo}, \texttt{c1}, \texttt{vsum}, \texttt{thl},
\texttt{sumj}. Its physical-space reference is Baez~\cite{Baez},
the section ``Gauge Theory on a Graph,'' Lemmas 1, 2 and the source
lemma \texttt{lem2.5}. Baez's link convention
$A_e\mapsto h_{t(e)}A_eh_{s(e)}^{-1}$ is related to the input's
$U_e\mapsto h_{s(e)}U_eh_{t(e)}^{-1}$ by the involution
$\jmath(U)_e=U_e^{-1}$: direct multiplication gives
$\jmath(\alpha_hU)_e=h_{t(e)}\jmath(U)_eh_{s(e)}^{-1}$.
The pullback is a Haar-unitary map with itself as inverse,
preserves $H^1,H^2$ because inversion is a bi-invariant metric
isometry, and intertwines each $E_e$. The raw Wilson products in
our Hamiltonian are unchanged. Both original TeX sources are retained.

The source conditional complement retains individual fine links.
Our complement retains their ordered products. The following proves
the exact relation, including the Hilbert and operator domains.

\begin{proposition}\label{prop:physicalbridge}
Let $2\leq m\mid2L$, $s=(U_e)_{e\in S_m}$,
$r=(U_e)_{e\in R_m}$, and define
\[
 \mathcal A_m=\SU(2)^{R_m},\quad
 \bar\rho_R(r)=\int\rho(s,r)\dd\lambda_{S_m}(s),\quad
 \mathcal K_R=L^2(\mathcal A_m,\bar\rho_R\lambda_R).
\]
Let $J_S:\mathcal K_R\to\HH$ be $J_SF(s,r)=F(r)$, with
\[
 (J_S^*u)(r)=\frac{1}{\bar\rho_R(r)}
                \int\rho(s,r)u(s,r)\dd\lambda_{S_m}(s),
 \qquad P_S=J_SJ_S^*.
\]
Write $\mathcal K_R^{\rm G}$ for invariance under the endpoint action
of all original vertices. There is a unitary map of physical spaces
\begin{equation}\label{eq:skeletonunitary}
 \iota_m:\KK_m^{\rm G}\longrightarrow\mathcal K_R^{\rm G},
 \qquad (\iota_m f)(r)=f((u_{\gamma,1}\cdots u_{\gamma,m})_\gamma).
\end{equation}
Its inverse on arbitrary $L^2$ classes is
\begin{equation}\label{eq:skeletoninverse}
 (\iota_m^{-1}F)(z)=\int F(\Theta_m^{-1}(v,z))\dd\lambda_v(v),
 \quad
 \Theta_m:\mathcal A_m\longrightarrow
 \left(\prod_\gamma\SU(2)^{m-1}\right)\times\mathcal Q_m,
 \quad r\longmapsto(v,z).
\end{equation}
The inverse coordinates are $u_{\gamma,k}=v_{\gamma,k}$ for $k<m$
and $u_{\gamma,m}=(v_{\gamma,1}\cdots v_{\gamma,m-1})^{-1}z_\gamma$.
Both maps identify the invariant $H^1$ and $H^2$ spaces, and
\begin{equation}\label{eq:physicalprojections}
 J_S\iota_m=J_m,
 \qquad J_S^*|_{\HH^{\rm G}}=\iota_mJ_m^*|_{\HH^{\rm G}},
 \qquad P_S|_{\HH^{\rm G}}=P_m|_{\HH^{\rm G}}.
\end{equation}
Consequently the identity on functions is an isometry from the
source physical conditional complement for $S=S_m$ onto
$\HH_{Q_m}^{\rm G}$. It identifies their full form domains,
their complementary operator domains, their operators and their
resolvents. In the range \eqref{eq:stronggap},
\begin{equation}\label{eq:uniformC}
 C\geq\Delta_{\rm sc} I,\qquad
 \|C^{-1}\|\leq\Delta_{\rm sc}^{-1}
                  \leq\frac{96}{287\kappa}.
\end{equation}
\end{proposition}

\begin{proof}
The two inverse coordinate compositions and the Haar identity for
$\Theta_m$ follow by the exact prefix cancellations already proved
for $\Phi_m$, with all $S_m$ factors omitted from both sides of
this particular map. Integrating those factors in the density gives
\[
 \bar\rho_R(\Theta_m^{-1}(v,z))=r_m(z),\qquad
 \bar\rho_R\dd\lambda_R=r_m(z)\dd\lambda_v\dd\lambda_m.
\]
To prove the first equality, perform simultaneously the internal
vertex transformations with fixed coarse endpoints as in Section 3;
they send $(v,z)$ to $(1,z)$, preserve the marginal by substitution
on $S_m$, and leave $z$ unchanged. Integrating in $v$ then identifies
the remaining function with the definition of $r_m$.

There is also an $L^2$ proof of the descent, without evaluation on
a measure-zero section. For a chain with endpoints fixed, the
internal gauge action is
$v'_1=u_1h_1$ and $v'_k=h_{k-1}^{-1}u_kh_k$ for $1<k<m$.
For fixed input $r$, Haar integration successively over
$h_1,\ldots,h_{m-1}$ makes the variables $v'_1,\ldots,v'_{m-1}$
product Haar; the chain product remains $z$. Applying this to every
disjoint interior shows that the compact internal gauge average on
$L^2(\bar\rho_R\lambda_R)$ is precisely fibre integration in $v$
followed by constant lift. Hence every invariant $F$ equals that
lift almost everywhere. The remaining coarse endpoint action is
$z_\gamma\mapsto h_x^{-1}z_\gamma h_{x+m\mathbf e_i}$; this proves
the full coarse invariance of \eqref{eq:skeletoninverse} and the
fine invariance of \eqref{eq:skeletonunitary}. Substitution proves
both inverse identities. Moreover
\[
 \|\iota_m f\|_{\mathcal K_R}^2
 =\int |f(z)|^2r_m(z)\dd\lambda_v\dd\lambda_m
 =\|f\|_{\KK_m}^2.
\]
No vector is divided by its mass. Smooth coordinate pullback on
these compact products preserves each Sobolev order. Differentiating
under the compact fibre integral and applying Cauchy--Schwarz bounds
the inverse in $H^1$ and $H^2$. Smooth positive weights give the
same Sobolev sets in each finite box. Approximation, followed by
compact gauge averaging, proves the claimed domain identifications.

The first identity in \eqref{eq:physicalprojections} follows from
the coordinate product. For physical $u$, substitution in the
conditional integral shows that $J_S^*u$ is invariant under every
skeleton endpoint action. In $(v,z)$ coordinates it is therefore
independent of $v$ almost everywhere. Its $v$ integral equals
$r_m(z)^{-1}\int\rho(s,\Theta_m^{-1}(v,z))
u(s,\Theta_m^{-1}(v,z))\dd\lambda_{S_m}\dd\lambda_v$,
which is exactly $J_m^*u$. This proves the second identity and then
the third. It does not assert equality of the projections on all
noninvariant vectors.

Thus the two physical complements are the same closed subspace
of $\HH$, with the same form $q$, including derivatives in every
fine link. Both have form domain $H^1$ intersected with that space.
The weak operator-domain rule is, in both presentations,
\[
\begin{aligned}
 D(C)=\bigl\{u\in H^1\cap\HH_{Q_m}^{\rm G}:{}&
 \text{ there exists }w\in\HH_{Q_m}^{\rm G}\text{ such that}\\
 &q(v,u)=\langle v,w\rangle
 \text{ for every }v\in H^1\cap\HH_{Q_m}^{\rm G}\bigr\},\\
 Cu={}&w.
\end{aligned}
\]
Equality of the tests, inner products and forms proves equality
of the operators and their resolvents, not merely of their formal
differential expressions. Finally every vector in the complement
has $\mu u=0$, and $\mathcal Uu=\psi u$ is physical and perpendicular
to $\psi$. The already-proved \eqref{eq:stronggap} applied under
this unchanged unitary gives \eqref{eq:uniformC}.
\end{proof}

\begin{theorem}\label{thm:uniformzero}
In the range \eqref{eq:stronggap}, for every mean-zero
$f\in\VV_m\cap\KK_m^{\rm G}$, put $b_f=B_mf$,
$d_f=\|f\|^2$, $a_f=a_m(f,f)$ and retain the raw measure
$\sigma_f(I)=\|1_I(C)b_f\|^2$.
Then
\begin{equation}\label{eq:uniformzeroenergy}
 s_{m,0}(f,f)\geq\Delta_{\rm sc}
              (d_f+\|C^{-1}b_f\|^2),\qquad
 \lambda_m^{\rm R}\geq\Delta_{\rm sc},\quad c\geq\Delta_{\rm sc}.
\end{equation}
For the unchanged mean-centered simple loop $f_C$, let
$\ell$ be its fine perimeter and $d_C,r_C$ retain their earlier
definitions. The physical zero-shift state $u_0=\psi\mathcal R_m f_C$
has the exact mass and energy
\begin{align}
 \|u_0\|^2&=d_C+\int t^{-2}\dd\sigma_{f_C}(t),\label{eq:uniformloopmass}\\
 \mathfrak q_{H-E_0}[u_0]
 &=\kappa\ell(1-r_C/4)-\int t^{-1}\dd\sigma_{f_C}(t)
 \geq\Delta_{\rm sc}\|u_0\|^2.\label{eq:uniformloopenergy}
\end{align}
Its quantitative upper mass bound is
\begin{equation}\label{eq:uniformloopbound}
 \|u_0\|^2\leq d_C+
 \min\left\{\frac{\kappa\ell(1-r_C/4)}{\Delta_{\rm sc}},
             \frac{16G^2\kappa^2\ell^2}{\Delta_{\rm sc}^2}\right\}
 \leq4+\frac{96\ell}{287}.
\end{equation}
For every $z\geq0$, define $w_z=(C+z)^{-1}b_f$,
$N_f(z)=d_f+\|w_z\|^2$ and
$e_f(z)=a_f-\int(t+z)^{-1}\dd\sigma_f
                    -z\int(t+z)^{-2}\dd\sigma_f$.
These are the full physical mass and excitation energy of
$\psi(J_mf-w_z)$. The exact differences are
\begin{align}
 w_0-w_z&=zC^{-1}(C+z)^{-1}b_f,\label{eq:uniformwdiff}\\
 e_f(z)-e_f(0)&=\int\frac{z^2}{t(t+z)^2}\dd\sigma_f(t),
                                                        \label{eq:uniformediff}\\
 N_f(0)-N_f(z)&=\int\left(\frac1{t^2}-\frac1{(t+z)^2}\right)
                                      \dd\sigma_f(t),    \label{eq:uniformndiff}\\
 s_{m,z}(f,f)-s_{m,0}(f,f)-zd_f
 &=\int\frac{z}{t(t+z)}\dd\sigma_f(t).    \label{eq:uniformmemorydiff}
\end{align}
Writing $\theta=z/(\Delta_{\rm sc}+z)$, their estimates are
\begin{align}
 \|w_0-w_z\|&\leq\theta\sqrt{a_f/\Delta_{\rm sc}},\nonumber\\
 0\leq e_f(z)-e_f(0)&\leq\theta^2a_f,\nonumber\\
 0\leq N_f(0)-N_f(z)&\leq
 \left[1-\left(\frac{\Delta_{\rm sc}}{\Delta_{\rm sc}+z}\right)^2\right]
 \frac{a_f}{\Delta_{\rm sc}},\nonumber\\
 0\leq s_{m,z}(f,f)-s_{m,0}(f,f)-zd_f&\leq\theta a_f.
                                                        \label{eq:uniformcontinuity}
\end{align}
The constants controlling the inverse do not depend on box or block
size. The displayed $a_f$, $\ell$, $G$, $\kappa$, raw norms and
spectral measures keep their full dependence.
\end{theorem}

\begin{proof}
The map $\mathcal R_m$ and its exact domain and energy identities
were proved in Proposition~\ref{prop:finitezero}. Its image under
$\mathcal U$ is physical and perpendicular to $\psi$. Applying
\eqref{eq:stronggap} proves \eqref{eq:uniformzeroenergy}, including
the infimum inequality. The previous exact inverse-moment formulas
give \eqref{eq:uniformloopmass}--\eqref{eq:uniformloopenergy}.
Positivity of $s_{m,0}$ gives
$\int t^{-1}\dd\sigma_f\leq a_f$. Since
$\operatorname{supp}\sigma_f\subset[\Delta_{\rm sc},\infty)$,
\[
 \int t^{-2}\dd\sigma_f\leq
 \min\{a_f/\Delta_{\rm sc},\|b_f\|^2/\Delta_{\rm sc}^2\}.
\]
Insert the original loop identities $a_f=\kappa\ell(1-r_C/4)$,
$\|b_f\|\leq4G\kappa\ell$, $d_C\leq4$, $a_f\leq\kappa\ell$
and \eqref{eq:stronggap} to prove \eqref{eq:uniformloopbound}.

The resolvent identity proves \eqref{eq:uniformwdiff}. Subtraction
of the complete energy integrands gives
$t^{-1}-(t+z)^{-1}-z(t+z)^{-2}=z^2/[t(t+z)^2]$.
The norm subtraction and the definition of $s_{m,z}$ give the
remaining exact differences. On the support of the raw measure,
$z/(t+z)\leq\theta$ and
$1-[t/(t+z)]^2\leq1-[\Delta_{\rm sc}/(\Delta_{\rm sc}+z)]^2$.
Apply these inequalities to the corresponding nonnegative
integrands and use the preceding inverse-moment bounds. Squaring
\eqref{eq:uniformwdiff} first proves its norm estimate. This proves
every inequality in \eqref{eq:uniformcontinuity}, with no removal
of the measure or its total mass $\|B_mf\|^2$.
\end{proof}

For an elementary coarse face $\ell=4m$, the last mass bound is
$4+384m/287$, and its physical side length remains $ma$.
In original parameters the new lower energy scale and range are
\[
 \Delta_{\rm sc}=\frac{2g^2}{a}\left(3-\frac{128}{g^4}\right)
 \geq\frac{287g^2}{48a},\qquad g^4\geq12288.
\]
The expectation of $H$ is still $E_0N_f(z)+e_f(z)$.
This is a statement for the interacting finite boxes, uniform as
their sizes vary within this explicit coupling range. It takes no
spatial continuum limit and provides no small-$g$ continuation.
Every induced score covariance and the full resolvent and time memory
remain those of Sections 4--6.
